\subsection{Locally path-connected spaces}

DEFINITION 4. --- A topological space is said to be locally path-connected if each of its points has a fundamental system of path-connected neighborhoods.

PROPOSITION 7. --- A locally path-connected topological space is locally connected. Its connected components coincide with its path-connected components. In particular, if a locally path-connected topological space is connected, it is path-connected.

The first assertion follows from Prop. 4. Let us prove the second assertion. Let X be a locally path-connected topological space and let C be a path-connected component of X. Every point of C has a path-connected neighborhood, hence contained in C; consequently, C is an open subset of X. Since the path-connected components form a partition of X, such a component is also closed. As it is connected (III, p. 258, prop. 4), it is a connected component (TG, I, p. 83). The third assertion follows from the second.

Note thus that there is no ambiguity in saying that a topological space is connected and locally path-connected.

COROLLARY 1. --- Every connected open subset of a locally path-connected topological space is path-connected.

COROLLARY 2. --- Let B be a locally path-connected topological space and let E be an étale B-space. The space E is locally path-connected. If it is connected, it is path-connected.

The first assertion follows immediately from def. 4 and the definition of an étale morphism (I, p. 28, def. 2). The second is deduced from it by prop. 7.

For a topological space X to be locally path-connected, it is necessary and sufficient that every path-connected component of an open subset of X be an open subset of X (cf. I, p. 85, proof of Prop. 11). If the space X is locally path-connected, then every point of X has a fundamental system of open path-connected neighborhoods.

PROPOSITION 8. --- Every quotient space of a locally path-connected space is locally path-connected.

Let X be a locally path-connected space, let R be an equivalence relation on X, and let \(\varphi\colon\mathrm{X}\to\mathrm{X}/\mathrm{R}\) be the canonical map. It suffices to show that the path-connected components of an open subset of X/R are open. Let A thus be an open subset of X/R and let C be a path-connected component of A. Let \(x\in\bar{\varphi}^{1}(\mathrm{C})\), and let K be the path-connected component of x in \(\bar{\varphi}^{1}(\mathrm{A})\). The set \(\varphi(\mathrm{K})\) is path-connected (III, p. 258, prop. 3), contained in A and contains \(\varphi(x)\); one therefore has \(\varphi(\mathrm{K})\subset\mathrm{C}\), whence \(\mathrm{K}\subset\bar{\varphi}^{1}(\mathrm{C})\). This proves that \(\bar{\varphi}^{1}(\mathrm{C})\) is a union of path-connected components of \(\bar{\varphi}^{1}(\mathrm{A})\). Since X is locally path-connected and \(\bar{\varphi}^{1}(\mathrm{A})\) is open in X, \(\bar{\varphi}^{1}(\mathrm{C})\) is open in X. Consequently, C is open in X/R. This proves the proposition.

PROPOSITION 9. --- The product of a family of locally path-connected spaces, all but finitely many of which are connected, is locally path-connected.

Let \((\mathrm{X}_{j})_{j\in\mathrm{J}}\) be a family of locally path-connected spaces that are connected, except for finitely many of them. Let \(x=(x_{j})_{j\in\mathrm{J}}\) be a point of the product space \(X=\prod_{j\in J}X_{j}\). By hypothesis, a fundamental system of neighborhoods of x in X consists of the sets of the form \(V=\prod_{j\in J}V_{j}\) where, for every \(j\in J\), \(V_{j}\) is a path-connected neighborhood of \(x_{j}\) in \(X_{j}\) and where \(V_{j}=X_{j}\) except for a finite set of indices \(j\in J\) (TG, I, p. 24). These sets being path-connected (III, p. 260, prop. 6), X is locally path-connected.

COROLLARY. --- Every open subset of a numerical space is locally path-connected.

Every point of R has a base of neighborhoods consisting of intervals; consequently, R is locally path-connected. By Proposition 9, the numerical space \(R^{n}\) is locally path-connected, for every integer \(n \geqslant 1\). Furthermore, an open subset of a locally path-connected topological space is still locally path-connected, whence the corollary.

Example. --- Let G be the subgroup of \(\mathbf{GL}(n,\mathbf{R})\) consisting of square matrices of order n whose determinant is strictly positive. The group G is connected and locally path-connected.

By A, III, p. 104, prop. 17, the group \(\mathbf{SL}(n,\mathbf{R})\) is generated by the elements \(B_{ij}(\lambda)\) (for \(1\leqslant i,j\leqslant n\) such that \(i\neq j\) and \(\lambda\in\mathbf{R}\)). The maps \(\lambda\mapsto B_{ij}(\lambda)\) from \(\mathbf{R}\) into \(\mathbf{GL}(n,\mathbf{R})\) are continuous, their images are connected subsets of \(\mathbf{SL}(n,\mathbf{R})\); since they all contain the identity matrix \(I_n\) their union is connected (TG, I, p. 81, prop. 2). Consequently, the group \(\mathbf{SL}(n,\mathbf{R})\) is connected (TG, III, p. 8, prop. 7). Let A be the subgroup of G formed by the matrices of the form diag(1, \(\ldots\),1, \(\lambda\)), with \(\lambda\in\mathbf{R}_{+}^{*}\); it is connected and one has \(\mathbf{GL}(n,\mathbf{R})=\mathbf{A}\cdot\mathbf{SL}(n,\mathbf{R})\). It follows that the group G is connected. Since it is the inverse image of \(\mathbf{R}_{+}^{*}\) under the determinant map from \(\mathbf{M}_{n}(\mathbf{R})\) into \(\mathbf{R}\), it is an open subset of \(\mathbf{M}_{n}(\mathbf{R})\); it is therefore locally path-connected (corollary above), as well as path-connected (III, p. 261, corollary 1).

PROPOSITION 10. --- Let X be a topological space. The map \((o, e)\) from \(\Lambda(\mathrm{X}) = \mathcal{C}_{c}(\mathbf{I}; \mathrm{X})\) into \(X \times X\), which associates to a path c in X the pair \((c(0), c(1))\), is continuous. If the space X is locally path-connected, this map is open.

We already know that the origin map o and the endpoint map e of \(\Lambda(\mathrm{X})\) into X are continuous (III, p. 257), whence the first assertion.

Lemma. --- Let \(c\colon \mathbf{I}\to \mathrm{X}\) be a path and let W be a neighborhood of \(c\) in the space \(\mathcal{C}_c(\mathbf{I};\mathrm{X})\). There exists a real number \(\varepsilon \in ]0,1 / 2]\), a neighborhood \(\mathrm{V}_0\) of \(c(0)\) and a neighbourhood \(\mathrm{V}_1\) of \(c(1)\) in X having the following properties: we have \(c([0,\varepsilon])\subset \mathrm{V}_0\), \(c([1 - \varepsilon,1])\subset \mathrm{V}_1\), and every path \(c^{\prime}\colon \mathbf{I}\rightarrow \mathrm{X}\) such that

\[
\left\{ \begin{array}{l l} c ^ {\prime} (t) \in \mathrm{V} _ {0} & \text { for } 0 \leqslant t \leqslant \varepsilon , \\ c ^ {\prime} (t) = c (t) & \text { for } \varepsilon \leqslant t \leqslant 1 - \varepsilon , \\ c ^ {\prime} (t) \in \mathrm{V} _ {1} & \text { for } 1 - \varepsilon \leqslant t \leqslant 1, \end{array} \right.\tag{2}
\]

belongs to W.

By definition of the topology of compact convergence (TG, X, p. 26, def. 1), there exists a finite set J, a family \((\mathrm{U}_{j})_{j\in\mathrm{J}}\) of open sets in X and a family \((\mathrm{K}_{j})_{j\in\mathrm{J}}\) of compact subsets of I such that the set \(W'\) of paths \(c'\) satisfying \(c'(\mathrm{K}_{j})\subset\mathrm{U}_{j}\) for every index j is a neighborhood of c contained in W. Let us then denote by \(A_{0}\) (resp. \(A_{1}\) ) the set of indices j such that \(0\in K_{j}\) (resp. \(1\in K_{j}\) ); set \(V_{0}=\bigcap_{j\in A_{0}}U_{j}\) and \(V_{1}=\bigcap_{j\in A_{1}}U_{j}\) .

Since the map \(c\) is continuous, there exists a real number \(\varepsilon \in [0,1/2]\) such that \(c([0,\varepsilon]) \subset \mathrm{V}_0\) , \(c([1-\varepsilon,1]) \subset \mathrm{V}_1\) , \([0,\varepsilon] \cap \mathrm{K}_j = \emptyset\) for all \(j \notin \mathrm{A}_0\) and \([1-\varepsilon,1] \cap \mathrm{K}_j = \emptyset\) for all \(j \notin \mathrm{A}_1\) Then let \(c'\) be a path satisfying conditions (2). Let us prove that \(c' \in \mathrm{W}'\) . Let \(j \in \mathrm{J}\) and let \(t \in \mathrm{K}_j\) . If \(\varepsilon \leqslant t \leqslant 1-\varepsilon\) , \(c'(t) = c(t)\) belongs to \(\mathrm{U}_j\) . If \(0 \leqslant t \leqslant \varepsilon\) , \(c'(t) \in \mathrm{V}_0\) ; by the choice of \(\varepsilon\) , we have \(j \in \mathrm{A}_0\) , therefore \(c'(t) \in \mathrm{U}_j\) . Similarly, if \(1-\varepsilon \leqslant t \leqslant 1\) , we have \(j \in \mathrm{A}_1\) and \(c'(t) \in \mathrm{V}_1 \subset \mathrm{U}_j\) . Thus, \(c'(\mathrm{K}_j) \subset \mathrm{U}_j\) and \(c'\) belongs to \(\mathrm{W}'\) , hence to W.

Let us now prove the second assertion of prop. 10. Suppose the space X is locally path-connected. Let \(c\) be a path in X and let W be a neighborhood of \(c\) in \(\mathcal{C}_c(\mathbf{I};\mathbf{X})\). Let \(\varepsilon\), \(\mathrm{V}_0\), and \(\mathrm{V}_1\) be as in the lemma. Let \(\mathrm{T}_0\) and \(\mathrm{T}_1\) be path-connected neighborhoods of \(c(0)\) and \(c(1)\) contained in \(\mathrm{V}_0\) and \(\mathrm{V}_1\), respectively. There exists a real number \(\theta\) such that \(0 < \theta < \varepsilon\) and such that \(c([0,\theta]) \subset \mathrm{T}_0\), \(c([1 - \theta,1]) \subset \mathrm{T}_1\). Let \(x_0 \in \mathrm{T}_0\) and \(x_1 \in \mathrm{T}_1\); let \(c_0\) be a path of origin \(x_0\) and of endpoint \(c(\theta)\) in \(\mathrm{T}_0\) and let \(c_1\) be a path of origin \(x_1\) and of endpoint \(c(1 - \theta)\)

in \(\mathrm{T}_1\). Set

\[
c ^ {\prime} (t) = \left\{ \begin{array}{l l} c _ {0} (t / \theta) & \text { for } 0 \leqslant t \leqslant \theta , \\ c (t) & \text { for } \theta \leqslant t \leqslant 1 - \theta , \\ c _ {1} ((1 - t) / \theta) & \text { for } 1 - \theta \leqslant t \leqslant 1. \end{array} \right.
\]

One thus defines a path \(c'\) connecting \(x_{0}\) to \(x_{1}\) and satisfying conditions (2). This proves that the image of W in X × X under the map \((o, e)\) contains the neighborhood \(T_{0} \times T_{1}\) of \((c(0), c(1))\) , whence the proposition.

COROLLARY. --- Let X be a locally path-connected topological space and let x be a point of X. The map \(c \mapsto c(1)\) of \(\Lambda_{x}(\mathrm{X})\) into X is open.

By the proposition, the map \(\varphi\colon\Lambda(\mathrm{X})\to\mathrm{X}\times\mathrm{X}\) defined by \(\varphi(c)=(c(0),c(1))\) is open and one has \(\Lambda_{x}(\mathrm{X})=\overline{\varphi}^{1}(\{x\}\times\mathrm{X})\) . Consequently, the map \(\Lambda_{x}(\mathrm{X})\to\{x\}\times\mathrm{X}\) deduced from \(\varphi\) is open (TG, I, p. 30, prop. 2), as is its composite with the second projection \(\mathrm{pr}_{2}\) .

